\uuid{JKUg}
\titre{Fonction croissante/décroissante}
\auteur{bodin ; hanani ; mzari}
\organisation{exo7}
\licence{CC BY-NC-SA 4.0}
\datecreate{2020-01-20}
\niveau{L1}
\module{Analyse}
\chapitre{Continuité, limite et étude de fonctions réelles}
\sousChapitre{Autre}
\difficulte{3}
\theme{Notion de fonctions, fonction croissante, décroissante}
\exo7qcmid{lille-1.continuite.008}

\contenu{
\begin{question}
\qclassification{123.00}
\qtags{motcle=fonction croissante/décroissante}
Soit $f : \Rr \to \Rr$ une fonction à valeurs strictement positives. Quels arguments sont valables pour montrer que $f$ est croissante ?
\begin{answers}
    \bad{Pour tout $x\in\Rr$, on a $f(x+1) \ge f(x)$.}

    \bad{Pour tout $x\in\Rr$, on a $\frac{f(x+1)}{f(x)} \ge 1$.}

    \bad{Pour tout $x\in\Rr$, il existe $h>0$ tel que $f(x+h) \ge f(x)$.}

    \good{Pour tout $x\in\Rr$, pour tout $h>0$, on a $\frac{f(x+h)}{f(x)} \ge 1$.}

\end{answers}
\begin{explanations}
En notant $y = x+h$ avec $h>0$ on a $y > x$ et donc il suffit d'avoir $\frac{f(y)}{f(x)} \ge 1$.
Par contre, il n'est pas suffisant de comparer $f$ en des valeurs distantes de $1$ ! Essayez de dessiner un contre-exemple : $f$ vaut $0$ partout, sauf $1$ en chaque entier.

\end{explanations}
\end{question}
}
